% SHARD-ID: LB-07-MEMORY-INDEX
% SHARD-TITLE: The memory index: quantization without scattering theory
% SHARD-SOURCES: definitions.md D26--D27; claims/CLAIMS.md rows M-INDEX-fin,
%   M-INDEX-spec, M-IDX-density, M-INDEX-LA-folium, M-INDEX-LA-strong,
%   LR1-GEN, M-SCOPE-center, M-INDEX-torus (adjudicated scoping verbatim);
%   theory/memory-index.md; theory/folium-implementer.md;
%   theory/compact-g-memory-scope.md; theory/lr-d16.md (step (1.3) only).
% SHARD-CLAIMS: D26, D27, M-INDEX-fin, M-INDEX-spec, LR1-GEN, M-IDX-density,
%   M-INDEX-LA-folium, M-INDEX-LA-strong, M-SCOPE-center, M-INDEX-torus
%
% Local semantic macros (this shard only; they mirror notation.md rows that
% have no macro in main.tex's preamble).
\providecommand{\dephase}{\mathcal{D}}      % TPM dephasing (pinching) map
\providecommand{\specres}{E}                % spectral resolution of a charge
\providecommand{\tr}{\operatorname{tr}}     % matrix trace on the virtual space

\section{The memory index: quantization without scattering theory}
\label{sec:memory-index}

\subsection{The inversion of the continuum logic}

In the continuum infrared triangle the memory effect is the last of the three
corners to be quantized, and it is quantized from the outside. One first
builds a scattering theory: wave operators, an inventory of asymptotic
channels, and enough completeness to say that the inventory is exhaustive.
A soft theorem is then read off the $S$-matrix, and the memory --- the
permanent displacement left behind after everything has dispersed --- inherits
whatever discreteness the channel labels happen to carry. Every step of that
chain is analytically expensive, and on a lattice model of any interest not
one of the steps is available.

On the lattice the logic runs the other way. The quantum of memory is fixed
\emph{before} any dynamics is discussed, by two ingredients that a lattice
model supplies for free:

\begin{enumerate}
\item \emph{An integrality hypothesis on the on-site charge.} The selected
  charge generates a circle, so its $2\pi$ rotation acts on one site by a
  scalar. That single algebraic fact --- hypothesis (INT), Definition
  \ref{def:int} --- pins the on-site spectrum into one coset of $\Z$.
\item \emph{Window arithmetic.} The regularised charge of a finite window is a
  sum of on-site charges plus a real background number. Its spectrum
  therefore also lies in one coset of $\Z$; the coset is some fixed real
  offset $\kappa_{W,c_0}$ that depends on the window and the cut, and on
  nothing else.
\end{enumerate}

The measurement protocol does the rest. If one measures the \emph{same}
windowed observable twice --- once early, once late --- both readouts are drawn
from that same coset, and their difference is an integer no matter how
complicated the intervening evolution was, and no matter that the two
Heisenberg observables fail to commute. The offset does not have to be
computed; it has to cancel, and it cancels because it is the same offset
twice. Time evolution enters only through the fact that it is a
$*$-automorphism, so that it preserves spectra.

That is the whole mechanism, and it costs no wave operator, no channel
inventory, and no completeness. What scattering theory buys, when it happens
to be available, is exactly the residual question: \emph{which} integer, and
with what probability. Corollary \ref{thm:reduction-had} makes that precise
--- under the scattering hypotheses of \cref{sec:wave-operators} the law
collapses to a two-point law and reproduces the channel result of
\cref{sec:memory-quantization} --- but the support statement was already true
without it.

Two disciplines are binding throughout and are worth stating before anything
else, because most of the ways of misreading this section violate one of them.

\begin{remark}[Only outcomes are quantized, never their mean]
\label{ref:register-discipline}
The wall displacement $\delta x$ is an average over a probability law. It is
never claimed to be quantized, and generically it is not. Quantization
attaches to two objects only: the individual outcomes of the explicit
two-measurement protocol below, and --- separately, under the scattering
hypotheses --- the spectrum of the single displacement operator $\Delta X$ of
\cref{sec:memory-quantization}. In particular, writing $\delta x$ as an
expectation over some coupling of two marginal laws is bookkeeping, not an
outcome law, since no such joint law is constructed.
\end{remark}

\begin{remark}[Differences of lattice-valued observables are not
lattice-valued]
\label{ref:noncommuting-landmine}
It is tempting to argue that both $\Qhat(t_-)$ and $\Qhat(t_+)$ have integer
spectrum, hence so does their difference. That is false. On $\C^2$ take
$Q_-=\ket{0}\bra{0}$ and $Q_+=\ket{+}\bra{+}$: both spectra are
$\{0,1\}\subset\Z$, yet
\begin{equation}
  \spec(Q_+-Q_-)=\Bigl\{-\tfrac{1}{\sqrt2},+\tfrac{1}{\sqrt2}\Bigr\}.
  \label{eq:noncommuting-landmine}
\end{equation}
No argument below ever forms such a difference. Integrality enters only
through the two \emph{spectral resolutions} of the same window observable, at
step (1.5).(2.3) of the proof shard --- that is, through readouts, not through
operator arithmetic.
\end{remark}

\provenance{register discipline L3; landmine exhibit L1}
{\texttt{theory/memory-index.md} \S0 (register paragraph) and \S6 (L1, L3)}
{\texttt{theory/checks/memory\_index\_check.py} IDX-C2 (the naive unmeasured
 difference must fail at $>10^{-6}$: the dephasing defect is generically
 nonzero)}
{\texttt{memory-index-r1.md} V3; \texttt{memory-index-r2.md}}

\subsection{The two definitions}

Throughout this section the setting is that of \cref{sec:memory-quantization}:
a compact on-site symmetry group with a covariant family of injective
matrix-product vacua, a distinguished kink pair $(\alpha,\beta)$, one selected
primitive circle direction $\xi$ common to both unbroken subgroups, and the
Hermitian on-site charge $S^z:=-\mathrm{i}\,q(\xi)$ built from it. The vacuum
densities are $\omega_\alpha(S^z)=+s$ and $\omega_\beta(S^z)=-s$ with $s>0$,
and $\wallX_W$ is the windowed wall-position observable of the frozen
definition recalled in \cref{sec:mps-setting}.

\begin{definition}[Circle-integral on-site charge, hypothesis (INT)]
\label{def:int}
For the selected Hermitian circle charge $S^z=-\mathrm{i}\,q(\xi)$,
hypothesis \textbf{(INT)} says that there is $c\in U(1)$ with
\begin{equation}
  e^{2\pi\mathrm{i}S^z_x}=c\,\mathbb{1}
  \label{eq:int}
\end{equation}
on one site. Equivalently, choosing the unique $\kappa\in[0,1)$ with
$c=e^{2\pi\mathrm{i}\kappa}$,
\begin{equation}
  \spec S^z\subset\kappa+\Z .
  \label{eq:int-coset}
\end{equation}
This holds for every spin-$S$ chain, with $c=e^{2\pi\mathrm{i}S}$.

No arithmetic condition on the real tail-density parameter $\rho$ is part of
(INT). In particular $\rho=0$ is admitted --- that is the AKLT cross-check
--- and the condition $s>0$ that calibrates the wall coordinate is a separate
hypothesis belonging to the kink setting, not to (INT). Under the additional
smoothness hypothesis (S) at both tails the density is instead \emph{derived}
to satisfy $e^{2\pi\mathrm{i}\rho}=c=e^{-2\pi\mathrm{i}\rho}$, hence
$2\rho\in\Z$; that is Theorem \ref{thm:density} below, a conclusion and not a
hypothesis.
\end{definition}

\provenance{D26}
{\texttt{definitions.md}, ``Memory-index campaign --- D26--D27'' (merged from
 \texttt{theory/memory-index.md} \S8.1)}
{\texttt{theory/checks/memory\_index\_check.py} IDX-C1 (finite-window offset
 and TPM law); \texttt{theory/checks/memory\_index\_probe.py} P2}
{\texttt{memory-index-r2.md} (objection 7 repair: $\rho=0$ admitted);
 \texttt{memory-index-r3.md}}

The second definition packages what has to be assumed about the long-time
behaviour of the charge history. It is the only place in the section where a
genuinely dynamical hypothesis is made, and it is assumed, never derived.

\begin{definition}[Charge-history local relaxation, hypothesis (LR)]
\label{def:lr}
Fix a kink-class vector $\Psi$ of the $\ell^1$ class of
\cref{sec:mps-setting}, a cut $c_0$, and a padded exhaustion
$W_m=[a_m,b_m]\uparrow\Z$ with $c_0\in W_m$ for all $m$. Put
\begin{equation}
  \Qhat_{W,c_0}:=2s\,(\wallX_W-c_0)\in\alg_W ,
  \label{eq:qhat-def}
\end{equation}
and let $\specres_{W,t}$ be the spectral resolution of its Heisenberg
translate $\Qhat_{W,c_0}(t)=\alpha_t(\Qhat_{W,c_0})$. Hypothesis
\textbf{(LR)} consists of the following three clauses.

\medskip\noindent\textbf{(LR1) Common-sequence Cesàro convergence.} There is
\emph{one} sequence $T_n\to\infty$ such that, for every fixed window $W$, the
two one-time Cesàro states
\begin{equation}
  \omega^{+}_{W,n}(A)=\frac{1}{T_n}\int_{T_n}^{2T_n}
     \braket{\Psi,\alpha_t(A)\Psi}\,\mathrm{d}t,
  \qquad
  \omega^{-}_{W,n}(A)=\frac{1}{T_n}\int_{-2T_n}^{-T_n}
     \braket{\Psi,\alpha_t(A)\Psi}\,\mathrm{d}t,
  \qquad A\in\alg_W ,
  \label{eq:lr1-states}
\end{equation}
and the double-Cesàro two-measurement laws
\begin{equation}
  p_{W,n}(\nu)=\frac{1}{T_n^{2}}\int_{T_n}^{2T_n}\!\mathrm{d}t_+
   \int_{-2T_n}^{-T_n}\!\mathrm{d}t_-
   \sum_{q}\bigl\lVert \specres_{W,t_+}(\{q-\nu\})\,
                       \specres_{W,t_-}(\{q\})\,\Psi\bigr\rVert^{2}
  \label{eq:lr1-laws}
\end{equation}
converge. The sum runs over $q\in\spec\Qhat_{W,c_0}$, and spectral values
that are absent contribute zero.

\medskip\noindent\textbf{(LR2) Vanishing first-moment measurement
back-action.} With the dephasing (pinching) map
$\dephase_{W,t_-}(A)=\sum_q \specres_{W,t_-}(\{q\})\,A\,
\specres_{W,t_-}(\{q\})$, the double-Cesàro average of
\begin{equation}
  \bigl\langle\Psi,\bigl[\dephase_{W,t_-}\bigl(\Qhat_{W,c_0}(t_+)\bigr)
     -\Qhat_{W,c_0}(t_+)\bigr]\Psi\bigr\rangle
  \label{eq:lr2-defect}
\end{equation}
tends to zero for every fixed $W$. This is a first-moment nondemolition
condition only; asymptotic commutativity of the two Heisenberg observables is
\emph{not} assumed.

\medskip\noindent\textbf{(LR3) First-moment tightness of the spatial
family.} Writing $p_W$ for the fixed-window time limit supplied by (LR1), the
family $\{p_{W_m}\}$ is first-moment tight:
\begin{equation}
  \lim_{M\to\infty}\ \sup_m\ \sum_{\abs{\nu}>M}(1+\abs{\nu})\,p_{W_m}(\nu)=0 .
  \label{eq:lr3}
\end{equation}

\medskip\noindent\textbf{(LR-conv) Convenience normalisation, optional.}
$p_{W_m}$ converges weakly to a probability $p$. This clause buys only
uniqueness of the \emph{value} $\delta x$ below; support quantization is
subsequence-free without it, by Prokhorov's theorem on the closed set $\Z$.

\medskip\noindent\textbf{The ordered wall expectation (a definition, not a
hypothesis).} $\delta x$ is
\begin{equation}
  \delta x:=\lim_m\bigl[\omega^{+}_{W_m}(\wallX_{W_m})
                       -\omega^{-}_{W_m}(\wallX_{W_m})\bigr],
  \label{eq:deltax-def}
\end{equation}
taken along the full sequence when (LR-conv) holds and along an (LR3)
subsequence otherwise. Its existence is a corollary of (LR1)--(LR3): at every
fixed $m$ one has the exact ledger identity
\begin{equation}
  \sum_\nu \nu\,p_{W_m}(\nu)
   =-2s\bigl[\omega^{+}_{W_m}(\wallX_{W_m})
             -\omega^{-}_{W_m}(\wallX_{W_m})\bigr],
  \label{eq:ledger-fixed-m}
\end{equation}
and (LR3) makes the left-hand side converge, hence the right-hand side. This
$\delta x$ is the ordered asymptotic value of the finite-time wall observable;
the finite-time $\delta x$ is unchanged. The order of limits is binding: the
infinite-volume dynamics is formed first, the fixed-window time limits
second, and the spatial exhaustion last. No plane-wave or soft-momentum
interchange is included anywhere.
\end{definition}

\provenance{D27}
{\texttt{definitions.md}, ``Memory-index campaign --- D26--D27''; the
 existence corollary for $\delta x$ is \texttt{theory/memory-index.md} step
 (1.7).(2.2)}
{\texttt{theory/checks/memory\_index\_check.py} IDX-C4 (a tight family with
 non-convergent first moment is rejected: bare tightness is not enough);
 \texttt{theory/checks/lr\_d16\_check.py} C1(a),(b)}
{\texttt{memory-index-r1.md} (objection 1: the former clause LR4 was a
 theorem and is deleted; the interaction-range-collar clause was cited by no
 step and is deleted); \texttt{memory-index-r2.md}}

\begin{remark}[Why (LR3) carries the weight $(1+\abs{\nu})$]
\label{ref:why-first-moment-tight}
Bare tightness is not enough, and one-point expectation convergence can never
supply the missing control. Two exhibits make this concrete. First,
$\mu_W=(1-\tfrac1W)\delta_0+\tfrac{1}{2W}\delta_{W}+\tfrac{1}{2W}\delta_{-W}$
has pointwise limit $\delta_0$ and exact signed first moment $0$, yet
$\sum_{\abs{\nu}>R}\abs{\nu}\mu_W(\nu)\ge 1$ for every $R$. Second,
$p_m=(1-\tfrac1m)\delta_0+\tfrac1m\delta_m$ is tight with weak limit
$\delta_0$ but has first moment identically $1$. The $(1+\abs{\nu})$ weight
in \eqref{eq:lr3} is exactly what prevents both loss of mass and loss of the
first moment.
\end{remark}

\provenance{D27(LR3); landmine exhibit L2}
{\texttt{theory/memory-index.md} \S6 (L2)}
{\texttt{theory/checks/memory\_index\_check.py} IDX-C4}
{\texttt{memory-index-b-r1.md} objection 2; \texttt{memory-index-r1.md}}

\subsection{Finite windows: the protocol, and integrality by offset
cancellation}

\subsubsection*{The two-projective-measurement protocol}

Everything downstream is a statement about one explicitly specified
experiment, so it is worth stating the experiment before the theorem. Fix a
finite window $W=[a,b]$, a cut $c_0\in W$, and two times $t_-<t_+$. The
observable is the single windowed regularised charge $\Qhat_{W,c_0}$ of
\eqref{eq:qhat-def}; the same window, the same cut, and the same background
are used at both times.

At time $t_-$ one performs a projective measurement of
$\Qhat_{W,c_0}(t_-)$, obtaining an eigenvalue $q_-=q$ with probability
$\lVert \specres_{W,t_-}(\{q\})\Psi\rVert^{2}$ and leaving the (unnormalised)
post-measurement vector $\specres_{W,t_-}(\{q\})\Psi$. At time $t_+$ one
measures the same window observable again, now Heisenberg-translated to
$t_+$, obtaining $q_+$. The joint law of the pair, indexed by the
\emph{escaped charge} $\nu:=q_--q_+$, is
\begin{equation}
  p_{W;t_-,t_+}(\nu)=\sum_{q\in\kappa_{W,c_0}+\Z}
   \bigl\lVert \specres_{W,t_+}(\{q-\nu\})\,
               \specres_{W,t_-}(\{q\})\,\Psi\bigr\rVert^{2}.
  \label{eq:tpm-law}
\end{equation}
This is the two-projective-measurement (TPM) law. Nothing in
\eqref{eq:tpm-law} presumes that the two Heisenberg observables commute;
$p_{W;t_-,t_+}$ is a genuine outcome law of a stated procedure, not a
spectral property of some difference operator (Remark
\ref{ref:noncommuting-landmine}).

The sign convention is fixed once. The measured window-charge increment is
$\Delta Q_W:=q_+-q_-$, and the escaped charge is $\nu=-\Delta Q_W$; the
ledger of \cref{sec:memory-quantization} then reads
\begin{equation}
  2s\,\delta x=\lim \mathbb{E}[\Delta Q_W]=-\sum_\nu \nu\,p_\nu,
  \qquad\text{i.e.}\qquad
  \delta x=-\frac{1}{2s}\sum_\nu \nu\,p_\nu .
  \label{eq:ledger}
\end{equation}

The first moment of \eqref{eq:tpm-law} has an exact closed form which is what
clause (LR2) is designed to control:
\begin{equation}
  \sum_\nu \nu\,p_{W;t_-,t_+}(\nu)
   =\bigl\langle \Qhat_{W,c_0}(t_-)\bigr\rangle
    -\bigl\langle \dephase_{W,t_-}\bigl(\Qhat_{W,c_0}(t_+)\bigr)\bigr\rangle .
  \label{eq:tpm-mean}
\end{equation}
The second term is the \emph{dephased} final expectation, not the bare one;
the difference between the two --- the dephasing defect --- is generically
nonzero, and replacing \eqref{eq:tpm-mean} by the naive unmeasured difference
$\langle\Qhat(t_-)\rangle-\langle\Qhat(t_+)\rangle$ is a registered checker
mutation that fails.

\subsubsection*{The finite-window theorem}

\begin{theorem}[Windowed memory is quantized without scattering theory]
\label{thm:fin}
\statusProved\
Assume (INT) of Definition \ref{def:int}; assume the calibration of
\cref{sec:memory-quantization} fixing the wall-coordinate parameter $s>0$ to
be the tail density; and fix a finite window $W=[a,b]$ and a cut $c_0\in W$.
Then:
\begin{enumerate}
\item[(i)] the regularised wall charge has, expanded,
  \begin{equation}
    \Qhat_{W,c_0}=\sum_{x=a}^{b}S^z_x+s\,(a+b-1-2c_0)\,\mathbb{1},
    \label{eq:qhat-expanded}
  \end{equation}
  and its spectrum lies in one coset of $\Z$,
  $\spec\Qhat_{W,c_0}\subset\kappa_{W,c_0}+\Z$, with offset
  \begin{equation}
    \kappa_{W,c_0}\equiv \abs{W}\kappa+s\,(a+b-1-2c_0)
    \pmod{\Z},\qquad \abs{W}=b-a+1;
    \label{eq:coset}
  \end{equation}
\item[(ii)] the same coset occurs at every time, i.e.
  $\spec\Qhat_{W,c_0}(t)\subset\kappa_{W,c_0}+\Z$ for every $t$; and
\item[(iii)] the TPM law \eqref{eq:tpm-law} is a probability law whose escaped
  increment satisfies $\nu\in\Z$, with no assumption that
  $\Qhat_{W,c_0}(t_-)$ and $\Qhat_{W,c_0}(t_+)$ commute.
\end{enumerate}
\end{theorem}

\begin{scope}
The hypotheses are exactly (INT), the density calibration $s>0$, a finite
window $W=[a,b]$, and a cut $c_0\in W$; the offset is
\eqref{eq:coset} and is time-independent \emph{because} the dynamics acts by
$C^*$-automorphisms, which preserve spectra --- no further dynamical input is
used or available. Integrality of the TPM increment is offset cancellation at
fixed $W$, and it is \emph{not} spectral arithmetic for a difference of
noncommuting operators (Remark \ref{ref:noncommuting-landmine}). The
supporting numerical probe certifies $\spec\Qhat_W\subset\Z$ \emph{by
construction of the observable} in the exactly solvable instantiation of
\cref{sec:kink-model} --- there the on-site charge and the vacuum density are
both exactly $\pm\tfrac12$, so every eigenvalue is an exact machine integer
for any state, evolved or not. That is an arithmetic certificate, not a
dynamical one, and it says nothing about the outcome law. This theorem is a
finite additive statement about a $0$-form charge: it implies no displacement
law, no channel inventory, and no ordered limit.
\end{scope}

\emph{Idea of proof.} By (INT), $\spec S^z\subset\kappa+\Z$ on one site. The
window charge \eqref{eq:qhat-expanded} is a sum of $\abs{W}$ commuting on-site
charges plus a real multiple of the identity, so its spectrum is a sum of
finite spectra shifted by that number; this gives \eqref{eq:coset}. Spectra
are automorphism-invariant, which gives (ii). For (iii), positivity of
\eqref{eq:tpm-law} is manifest, and summing over $\nu$ and $q$ gives $1$
because each of the two spectral resolutions sums to the identity (apply
Parseval twice, sequentially). Finally both readouts $q_-$ and $q_+$ are
drawn from the \emph{same} coset $\kappa_{W,c_0}+\Z$, because both
measurements use the same $W$, the same $c_0$, and the same background; hence
$\nu=q_--q_+\in\Z$. The offset is removed here, at fixed $W$, before any
limit is taken --- which is also why it cannot smear when the window is later
enlarged.

\provenance{M-INDEX-fin; D13(a), D26}
{\texttt{theory/memory-index.md} steps (1.1) and (1.5); the offset
 cancellation is step (1.5).(2.3); the mean identity \eqref{eq:tpm-mean} is
 step (1.6)}
{\texttt{theory/checks/memory\_index\_check.py} IDX-C1, IDX-C2 (ten
 registered rows; green exit $0$, \texttt{--red} exit $1$, RED-OK $10$ of
 $10$); \texttt{theory/checks/memory\_index\_probe.py} P2 (arithmetic
 certificate only)}
{\texttt{memory-index-r2.md} \S4; gates cleared in
 \texttt{memory-index-r3.md} \S0 V3}

\subsection{The limit law}

The finite-window theorem says nothing about what happens as the window grows
and the times recede. That step is where a hypothesis has to be paid, and
(LR) is that hypothesis. What follows is a conditional implication and is
stated as one.

\begin{theorem}[The ordered escaped-charge law is supported on the integers]
\label{thm:spec}
\statusProvedConditional\
Assume the kink-and-calibration hypotheses (1)--(4) of
\cref{sec:memory-quantization} --- in particular no scattering hypothesis and
no channel inventory --- together with (INT) of Definition \ref{def:int} and
(LR1)--(LR3) of Definition \ref{def:lr} in their tightness-only form. Then
every subsequential ordered TPM escaped-charge law $p$ is a probability
supported on $\Z$, and along that subsequence
\begin{equation}
  \delta x=-\frac{1}{2s}\sum_{\nu\in\Z}\nu\,p_\nu .
  \label{eq:idx2}
\end{equation}
Support quantization is subsequence-free; the optional clause (LR-conv) makes
$p$ and the value $\delta x$ unique and does nothing else. Consequently the
individual outcomes of this explicit history protocol --- not their average
--- are multiples of $1/(2s)$ in wall-displacement units.
\end{theorem}

\begin{scope}
The theorem is proved only as the conditional implication just displayed:
the kink hypotheses, (INT) and (LR1)--(LR3) in tightness-only form imply that
every (LR3)-subsequential ordered TPM law is a probability on $\Z$ satisfying
\eqref{eq:idx2} \emph{along that subsequence}. Support quantization is
subsequence-free; (LR-conv) buys uniqueness of the value $\delta x$ and
nothing else. Hypothesis (LR) is assumed, never derived, and no implication
from the scattering hypotheses of \cref{sec:wave-operators} to (LR) is
claimed anywhere: the frozen words of the local-decay clause there fix no
window-marginal topology and supply no uniform integrability, so a future
bridge would have to add the reading of Corollary \ref{thm:reduction-had}
explicitly and prove (LR) from it. No sector-wide total-charge operator is
constructed here, and its unconditional existence is refuted --- see Refuted
claim \ref{ref:la-strong}. Bound states, absorption and additional channels
are all permitted, precisely because none of them is named.
\end{scope}

\emph{Idea of proof.} Three moves, in the order the limits are taken. First,
at fixed $W$, average the exact mean identity \eqref{eq:tpm-mean} doubly in
Cesàro form and pass to the common sequence of (LR1); clause (LR2) says
exactly that the dephasing defect disappears in that average, so what
survives is
\begin{equation}
  \sum_\nu \nu\,p_W(\nu)
   =\omega^{-}_{W}(\Qhat_{W,c_0})-\omega^{+}_{W}(\Qhat_{W,c_0}).
  \label{eq:spec-step1}
\end{equation}
This is the step that closes the two-time problem: it is (LR2), and not the
false arithmetic of Remark \ref{ref:noncommuting-landmine}, that does the
work. Second, substitute the definition \eqref{eq:qhat-def}; the common
scalar $-2sc_0$ cancels exactly and \eqref{eq:spec-step1} becomes the ledger
\eqref{eq:ledger-fixed-m}. Third, take the padded spatial exhaustion. Each
$p_{W_m}$ is a probability on the closed set $\Z$, so by Prokhorov every
subsequence has a further subsequence converging weakly to a probability
$p'$, necessarily supported on $\Z$ --- that part needs no more than
tightness, which is why it is subsequence-free. The $(1+\abs{\nu})$ weight in
(LR3) additionally passes the first moment to the limit, giving
\eqref{eq:idx2}. The window-dependent offset $\kappa_{W,c_0}$ cannot smear
under this limit because it was already removed at fixed $W$, in Theorem
\ref{thm:fin}(iii). No plane-wave or soft limit is taken: the normalisable
packet is fixed before all limits.

\provenance{M-INDEX-spec; depends on M-flux, M-quant-G, M-INDEX-fin, D13,
 D17, D26, D27}
{\texttt{theory/memory-index.md} steps (1.4)--(1.8); the Prokhorov step is
 (1.7).(2.3)}
{\texttt{theory/checks/memory\_index\_check.py} IDX-C1--IDX-C5 (ten
 registered rows; green exit $0$, \texttt{--red} exit $1$, RED-OK $10$ of
 $10$); \texttt{theory/checks/memory\_index\_probe.py} P1--P5 (PASS;
 \texttt{--red} exit $1$ with $190$ pre-registered violations,
 \texttt{--selftest} exit $0$)}
{\texttt{memory-index-r2.md} \S4; gates cleared in
 \texttt{memory-index-r3.md}}

\begin{corollary}[Reduction to the channel picture, when a channel picture
exists]
\label{thm:reduction-had}
\statusProvedConditional\
Assume in addition the scattering hypotheses of \cref{sec:wave-operators} in
the form used in \cref{sec:memory-quantization}, with the local-decay clause
read as \emph{weak-$*$ convergence of the window restriction, per channel, to
the corresponding kink charge eigenstate}. Then the limit law is the
two-point law
\begin{equation}
  p=(1-\braket{N_T})\,\delta_0+\braket{N_T}\,\delta_2 ,
  \label{eq:two-point}
\end{equation}
and \eqref{eq:idx2} together with its variance reproduces exactly the channel
result of \cref{sec:memory-quantization}:
\begin{equation}
  \delta x=-\frac{\braket{N_T}}{s},
  \qquad
  \frac{1}{(2s)^2}\operatorname{Var}_p(\nu)
    =\frac{1}{s^{2}}\braket{N_T}\bigl(1-\braket{N_T}\bigr).
  \label{eq:reduction}
\end{equation}
\end{corollary}

\begin{scope}
The reading of the local-decay clause is a further hypothesis, named here
because the frozen words of that clause fix no topology; it is the reading
the channel argument of \cref{sec:memory-quantization} needs for its own
uses. Reflection changes the core charge by $0$ and transmission by $-2$, so
the escaped values are $0$ and $2$; the two-point form \eqref{eq:two-point}
holds only inside the two-channel inventory, which excludes any further
propagating channel and any bound-state component of the selected vector.
Coherence is retained throughout: no replacement of the state by a norm
mixture is used.
\end{scope}

\emph{Idea of proof.} Two limits do the work. Incoming concentration: for the
incoming kink-charge eigenvalue $q_*$, the window projection
$\specres_W(\{q_*\})$ lies in $\alg_W$, and
$\lVert \specres_W(\{q_*\})\Psi_t-\Psi_t\rVert^{2}
 =1-\braket{\Psi_t,\specres_W(\{q_*\})\Psi_t}\to 0$ as $t\to-\infty$ by the
named reading. Cross-term vanishing: for the outgoing superposition
$r\ket{R}+t\ket{T}$, Cauchy--Schwarz gives
$\abs{\bra{R}\specres_W(\{q_R\})\ket{T}}
 \le \braket{T,\specres_W(\{q_R\})T}^{1/2}\to0$ since the transmitted leg's
window law concentrates on $q_T\neq q_R$. Hence the TPM Born weights receive
no inter-channel cross terms and are $1-\braket{N_T}$ and $\braket{N_T}$; the
rest is a Bernoulli computation.

\provenance{M-INDEX-spec (reduction clause)}
{\texttt{theory/memory-index.md} step (1.8), with the local-decay topology
 lemma at (1.8).(2.2)}
{\texttt{theory/checks/memory\_index\_check.py} IDX-C5 (mutating the
 transmitted escaped charge from $2$ to $1$ must fail both the mean and the
 variance)}
{\texttt{memory-index-r1.md} objection 4; \texttt{memory-index-r2.md};
 \texttt{memory-index-r3.md} objection 5 (indexing gloss withdrawn)}

\begin{remark}[An empirical fence for the channel-free register]
\label{ref:probe-fence}
The numerical probe of \cref{sec:numerics} finds, in the exactly solvable
chain of \cref{sec:kink-model}, extra outcome mass at window-charge values
$-1$ and $-3$, carried by near-threshold two-magnon channels with vanishing
group velocity. Those channels lie outside the two-channel inventory of
Corollary \ref{thm:reduction-had}, yet the mass sits exactly on the integers.
This is direct evidence for the channel-free register of Theorem
\ref{thm:spec}, and a concrete reason why the two-point form
\eqref{eq:two-point} needs its inelastic-threshold fence.
\end{remark}

\provenance{M-INDEX-spec (fence remark)}
{\texttt{theory/memory-index.md} \S4 remark, \S7 finding (i)}
{\texttt{theory/checks/memory\_index\_probe.py} P3 (the outcome-support
 finding is P3's, not P2's)}
{\texttt{memory-index-r2.md} objection 9; \texttt{memory-index-r3.md}}

\subsection{Clause (LR1) is a theorem}

Of the three clauses of (LR), the first is not a hypothesis at all. It
follows from soft-analysis inputs that any lattice model of the kind
considered here supplies. This matters twice over: it shortens what has to be
assumed, and it explains why the clause cannot simply be deleted.

\begin{theorem}[Cesàro compactness of the charge history]
\label{thm:lr1-gen}
\statusProved\
Let $\alg$ be the quasi-local algebra of a quantum spin system on a countable
lattice with uniformly finite local dimension, let $\alpha_t$ be a strongly
continuous one-parameter group of $*$-automorphisms of $\alg$, let $\Psi$ be a
unit vector in \emph{any} representation of $\alg$, let $c_0$ be a cut, and
let $\Qhat_{W,c_0}\in\alg_W$ be self-adjoint with finite spectrum at each
finite window $W\ni c_0$. Then \emph{every} prescribed sequence $S_n\to\infty$
admits a subsequence $T_n=S_{n_j}$ along which, \emph{simultaneously for every
such $W$}, the two Cesàro states $\omega^{\pm}_{W,n}$ of
\eqref{eq:lr1-states} converge weak-$*$ on all of $\alg$ to states, and every
double-Cesàro TPM weight $p_{W,n}(\nu)$ of \eqref{eq:lr1-laws} converges.
Each limit $p_W$ is a probability supported in the finite set
$\spec\Qhat_W-\spec\Qhat_W$, which under (INT) lies in $\Z$.
\end{theorem}

\begin{scope}
The proof consumes exactly three inputs: separability (and unitality) of
$\alg$; strong continuity of $\alpha_t$; and finiteness of
$\spec\Qhat_{W,c_0}$ at each fixed $W$. It uses \emph{no} spectral gap, no
Lieb--Robinson velocity, no ergodicity, no scattering or completeness input,
and no property of $\Psi$ beyond $\norm{\Psi}=1$. What is \emph{not} claimed:
pointwise-in-$t$ convergence, which (LR1) does not ask for; convergence along
the full sequence; and anything whatsoever about (LR2) or (LR3). One further
caveat is recorded: the ASSUME line of the proof step still over-lists two
standing register hypotheses that the argument does not use, and is due a
strengthening in a later round.

The clause is therefore a theorem rather than a restriction --- but it is not
removable. It existentially binds the \emph{single} sequence over which
clauses (LR2) and (LR3) are then quantified, and the substantive content of
(LR) lies entirely in those two clauses.
\end{scope}

\emph{Idea of proof.} The quasi-local algebra is a $C^*$-inductive limit over
a countable directed family of finite-dimensional matrix algebras, hence
separable: rational-coefficient words in fixed finite generating sets of the
local algebras form a countable norm-dense set. Therefore the state space
$S(\alg)$ is weak-$*$ compact \emph{and metrizable}, hence sequentially
compact, and two successive extractions produce one sequence along which both
$\omega^{+}$ and $\omega^{-}$ converge --- one sequence, all windows at once,
because convergence is on all of $\alg$. For the laws, fix $W$: since
$\Qhat_{W,c_0}$ has finite spectrum $\sigma_W$, its spectral projections lie
in $\alg_W$, the integrand of \eqref{eq:lr1-laws} is a finite sum of jointly
continuous functions of $(t_-,t_+)$ with values in $[0,1]$, it vanishes unless
$\nu\in\sigma_W-\sigma_W$, and it sums to $1$. The index set of pairs
$(W,\nu)$ is countable, so a Cantor diagonal extraction over
Bolzano--Weierstrass in $[0,1]$ yields one further sequence along which every
$p_{W,n}(\nu)$ converges. Finally a pointwise limit of probability measures
on a \emph{fixed finite} set is a probability measure on that set: no mass can
escape.

\provenance{LR1-GEN; D26, D27}
{\texttt{theory/lr-d16.md} steps (1.3).(2.1)--(2.8), restated at (1.6).(2.2);
 the scope inventory is step (1.3).(2.7)}
{\texttt{theory/checks/lr\_d16\_check.py} C1(a) TPM normalisation
 ($6.66\times10^{-16}$) and C1(b) raw integer spectrum and support; green
 exit $0$, \texttt{--red c1-nonunitary} and \texttt{--red c1-noninteger} each
 exit $1$ on exactly that row. These certify the finite-$W$ TPM arithmetic
 only; \emph{no} gate bears on the Cesàro convergence itself, which is not a
 finite-size testable statement and needs none --- the proof is complete}
{\texttt{lr-d16-r1.md}; \texttt{lr-d16-r2.md}}

\subsection{A by-product: the tail density is half-integer}

The circle hypothesis (INT) was introduced to control window spectra. Applied
to the two vacua themselves it yields something independent of any window, any
cut and any dynamics: a rigidity statement in the Lieb--Schultz--Mattis
family, proved here in the matrix-product register with no gap, no clustering
and no dynamical input beyond injectivity.

\begin{theorem}[Antisymmetric covariant tails have half-integer density]
\label{thm:density}
\statusProved\
Assume: injective canonical matrix-product tensors $A_\alpha,A_\beta$;
covariance of the vacuum family with a common unbroken circle direction
$\xi\in\mathfrak{h}_\alpha\cap\mathfrak{h}_\beta$ and on-site charge
$S^z=-\mathrm{i}\,q(\xi)$; the uniqueness clause of the intertwining relation;
the smoothness hypothesis (S) at both tails; (INT) of Definition \ref{def:int}
with scalar $c$; and the \textbf{antisymmetry}
$\omega_\beta(S^z)=-\omega_\alpha(S^z)=:-\rho$. Then the intertwining phase
$f_\alpha(\theta):=\theta_\alpha(\exp(\theta\xi))$ has slope
$f_\alpha'(0)=\rho$, satisfies $f_\alpha(\theta)=\rho\theta \bmod 2\pi$ for all
$\theta$, and
\begin{equation}
  e^{2\pi\mathrm{i}\rho}=c=e^{-2\pi\mathrm{i}\rho},
  \qquad\text{hence}\qquad e^{4\pi\mathrm{i}\rho}=1,
  \qquad 2\rho\in\Z,\qquad c\in\{\pm1\}.
  \label{eq:density-quantization}
\end{equation}
\end{theorem}

\begin{scope}
The antisymmetry is load-bearing. One tail alone gives only
$\rho\in\kappa+\Z$; a general (non-antisymmetric) tail pair gives only
$\rho_\alpha-\rho_\beta\in\Z$; the conclusion $2\rho\in\Z$ needs the
antisymmetric pair. The density $\rho$ is a free real parameter of the vacuum
pair, and $\rho=0$ is admitted --- that is the AKLT cross-check. No step
assumes any relation between $\rho$ and the on-site dimension $d$: $\rho$ is a
fresh density symbol and is deliberately \emph{not} the site-spin parameter
whose standard gloss $d=2s+1$ appears in the symbol table; that gloss is the
fully polarised special case and is not imported here. The result is derived,
not assumed.
\end{scope}

\emph{Idea of proof.} Four steps. (i) \emph{Slope equals density.} With
$V_\theta:=V_\alpha(\exp(\theta\xi))$ and $u(\theta):=e^{\mathrm{i}\theta S^z}$
on one site, the intertwining relation gives
$\omega_\alpha(u_x(\theta))
 =e^{\mathrm{i}f_\alpha(\theta)}\tr[V_\theta^{-1}E_\alpha(V_\theta r)]$.
Differentiating at $\theta=0$ --- legitimate by (S) --- the two virtual terms
cancel because the transfer operator fixes the right environment and preserves
traces, leaving
$\tfrac{\mathrm{d}}{\mathrm{d}\theta}\omega_\alpha(u_x(\theta))|_0
 =\mathrm{i}f_\alpha'(0)$, while the same derivative equals
$\mathrm{i}\,\omega_\alpha(S^z)=\mathrm{i}\rho$. Hence $f_\alpha'(0)=\rho$.
(ii) \emph{Additivity and continuity.} $f_\alpha$ is additive mod $2\pi$ along
the unbroken circle and is $C^1$ near the origin with $f_\alpha(0)=0$; the
additivity defect of a continuous lift is a continuous $2\pi\Z$-valued
function vanishing at the origin, hence identically zero nearby, and Cauchy's
functional equation gives $f_\alpha(\theta)=\rho\theta \bmod 2\pi$ globally.
(iii) \emph{Evaluate at $\theta=2\pi$.} (INT) gives $u(2\pi)=c\mathbb{1}$, so
the intertwining relation reads
$cA_\alpha^{\,s}=e^{\mathrm{i}f_\alpha(2\pi)}V_{2\pi}^{-1}A_\alpha^{\,s}V_{2\pi}$;
the uniqueness clause (applied with the trivial comparison gauge) forces
$V_{2\pi}$ scalar and $e^{\mathrm{i}f_\alpha(2\pi)}=c$, i.e.
$e^{2\pi\mathrm{i}\rho}=c$. (iv) \emph{The other tail.} The same computation
at $\beta$, whose density is $-\rho$ and whose on-site charge --- hence
whose scalar $c$ --- is the same, gives $e^{-2\pi\mathrm{i}\rho}=c$. Dividing
yields \eqref{eq:density-quantization}.

\provenance{M-IDX-density; D1, D2(a,b,e), D26}
{\texttt{theory/memory-index.md} step (1.9); route recorded in
 \texttt{memory-index-b-r1.md} \S1(b)}
{\texttt{theory/checks/memory\_index\_check.py} IDX-C6, IDX-C7(i),
 IDX-C7(ii) (green exit $0$, \texttt{--red} exit $1$). The load-bearing
 two-tail step (1.9).(2.4) is separately red-certified as IDX-C7(ii), with
 twisted-transfer defects at the correct counterterm and
 $\abs{\varphi_W(2\pi)-1}$ up to $1.9979$ on the charge-shifted family,
 independently re-run in \texttt{memory-index-r3.md} \S0 V3(i)--(iii)}
{\texttt{memory-index-b-r1.md}; \texttt{memory-index-r2.md} \S4 (objection 1
 fence, note 10); \texttt{memory-index-r3.md}}

\begin{remark}[The offset vanishes identically on an antisymmetric pair]
\label{ref:zero-offset}
If in addition the wall-coordinate calibration parameter is identified with
the tail density, $s=\rho$, then the coset offset \eqref{eq:coset} is not
merely constant in time but identically zero: substituting
$\kappa\equiv\rho \pmod{\Z}$ into \eqref{eq:coset} gives
\begin{equation}
  \kappa_{W,c_0}\equiv \rho(b-a+1)+\rho(a+b-1-2c_0)=2\rho\,(b-c_0)\equiv 0
  \pmod{\Z},
  \label{eq:zero-offset}
\end{equation}
since $2\rho\in\Z$ and $b-c_0\in\Z$. Hence $e^{2\pi\mathrm{i}\Qhat_{W,c_0}}$
is exactly the identity in $\alg_W$ and $\spec\Qhat_{W,c_0}\subseteq\Z$ for
every window and cut. Nothing in Theorems \ref{thm:fin}--\ref{thm:spec} needs
this: the TPM increment cancels the offset at fixed $W$ regardless. It merely
strengthens ``one coset'' to ``the zero coset'', and it is what the numerical
probe's integrality findings instantiate at $s=\tfrac12$.
\end{remark}

\provenance{M-IDX-density (corollary); D26, and the density calibration
 $s=\rho$}
{\texttt{theory/memory-index.md} step (1.10)}
{\texttt{theory/checks/memory\_index\_check.py} IDX-C7(i),(ii);\linebreak
 \texttt{theory/checks/memory\_index\_probe.py} P1, P2}
{\texttt{memory-index-r2.md} objection 4; \texttt{memory-index-r3.md}}

\subsection{An implementer on one kink folium}

Theorem \ref{thm:fin} is about finite windows. It is natural to ask whether
the window charges assemble into a single self-adjoint charge operator in an
infinite-volume representation. The next two results give the honest answer:
on one fixed kink representation such an operator exists, but it is
\emph{not} obtained as a limit of the window charges, and the general
existence statement is false.

\begin{theorem}[A strongly continuous circle implementer on the fixed kink
folium]
\label{thm:la-folium}
\statusProved\
Let $\varrho_0$ be the fixed bare kink state supplied by the corner-A
construction of \cref{sec:corner-a}, with invertible junction insertion
$V_\alpha(g)^{-1}$. Let the selected circle be common-unbroken,
$h_\theta\in H_\alpha\cap H_\beta$, satisfy the smoothness hypothesis (S) at
both tails, and obey (INT) of Definition \ref{def:int}. Then in the GNS
representation of $\varrho_0$ the circle automorphisms are implemented by a
strongly continuous unitary group
\begin{equation}
  U(\theta)=e^{\mathrm{i}\theta\Qhat},\qquad
  U(\theta)\pi(A)U(\theta)^{*}=\pi(\gamma_\theta(A)),\quad A\in\alg ,
  \label{eq:implementer}
\end{equation}
whose Stone generator satisfies
$e^{2\pi\mathrm{i}\Qhat}=e^{2\pi\mathrm{i}q_0}\mathbb{1}$ for some real
$q_0$; consequently $\Qhat$ is pure point with
$\spec\Qhat\subseteq q_0+\Z$, a single coset of $\Z$.
\end{theorem}

\begin{scope}
This is a folium-equivalence and implementer theorem. It is \emph{not} a
strong-resolvent limit of the window charges $\Qhat_{W,c_0}$: for a nonscalar
virtual circle action $V_\theta$ those bare window unitaries remain
non-Cauchy, as Refuted claim \ref{ref:la-strong} records, and there is no
contradiction because an implementing unitary is determined only up to a
scalar and is under no obligation to be the strong limit of any particular
inner approximants. The scope is the fixed bare kink of \cref{sec:corner-a},
together with only those finite-core variants whose padded blocked tensors
span the full matrix algebra $M_\chi(\C)$; no assertion is made for arbitrary
finite-core modifications. The common-unbroken recital is necessary: if the
circle moved either tail label, the rotated state would lie in a disjoint
sector and no implementer could exist in the fixed representation. The whole
content is \emph{existence}; once the implementer exists, the scalar-period
clause is a one-line Schur argument.
\end{scope}

\emph{Idea of proof.} Injectivity localises the circle orbit. Since the left
and right injectivity words each span $M_\chi(\C)$ and the junction matrix is
invertible, the blocked reference matrices $L_aMR_b$ span $M_\chi(\C)$; hence
a fixed padded block around the junction carries a local operator $B(M')$,
affine and continuous in $M'$ with $B(M)=\mathbb{1}$, such that
$\pi(B(M'))\Omega$ realises exactly the state with junction $M'$ replaced.
Rotating by the circle telescopes all virtual gauges away except at the
junction, where $M\mapsto M_\theta=V_\alpha(h_\theta)MV_\beta(h_\theta)^{-1}$;
this is continuous and invertible near $\theta=0$, so the rotated states are
represented by a continuous family of unit vectors $\xi_\theta$ in the
\emph{same} Hilbert space. The fixed kink is pure --- it is a vacuum composed
with a half-string automorphism --- so its GNS representation is irreducible,
and the rule
$W_\theta\pi(A)\Omega:=\pi(\gamma_\theta(A))\xi_{-\theta}$ defines an
isometry, extends to a unitary, and implements $\gamma_\theta$. Uniqueness of
implementers up to scalars makes $\theta\mapsto[W_\theta]$ a continuous
projective homomorphism, which on a \emph{one-dimensional} parameter group can
be de-projectivised explicitly: the $C^1$ scalar two-cocycle
$m(s,t)=W_sW_tW_{s+t}^{*}$ has argument $f(s,t)=a(s+t)-a(s)-a(t)$ for an
explicit $a$, so $e^{\mathrm{i}a(\theta)}W_\theta$ is an honest local group
and extends to $\R$. Stone's theorem produces $\Qhat$; (INT) makes
$\gamma_{2\pi}$ the identity automorphism, so $U(2\pi)$ lies in the scalar
commutant, and spectral calculus confines the spectral measure to
$q_0+\Z$ and makes it pure point.

\provenance{M-INDEX-LA-folium; A2, WI, D1(c,e$'$), D2(b,e), D26}
{\texttt{theory/folium-implementer.md} steps (1.1)--(1.7); the compatibility
 fence with Refuted claim \ref{ref:la-strong} is step (1.7); summarised at
 \texttt{theory/memory-index.md} step (1.3b)}
{--- (exact finite-dimensional and functional-analytic proof; no numerical
 claim and no checker required)}
{\texttt{blitz-la-folium-r1.md} (0 fatal / 0 major)}

\subsection{What is refuted: there is no sector-wide charge operator}

\begin{refutedclaim}[Unconditional existence of a regularised total charge]
\label{ref:la-strong}
\statusRefuted\
\emph{Withdrawn statement.} The kink-sector definition, the $\ell^1$
decay class, and (INT) together imply the existence of a self-adjoint
regularised total charge in every GNS representation of the kink sector
$\Kk_{\alpha\beta}$ --- the operator form of the finite-window theorem.

\emph{This is false, by two distinct mechanisms.}

\medskip\noindent\textbf{Mechanism A (off-folium; the fluctuation
counterexample).} This mechanism alone refutes the statement as written, and
it needs neither the smoothness hypothesis (S) nor any identification of the
calibration parameter with the density. Take one site with basis $\ket{m}$,
$m\in\{-\tfrac32,-\tfrac12,\tfrac12,\tfrac32\}$ and $S^z\ket{m}=m\ket{m}$; the
symmetry group $U(1)\rtimes\Z_2$ with the circle generated by $S^z$ and the
reflection $R\ket{m}=\ket{-m}$; the product vacua
$\alpha=\bigotimes\ket{\tfrac12}$ and $\beta=\bigotimes\ket{-\tfrac12}$; and
$H=0$. For $n\ge1$ put $\varepsilon_n=(n+1)^{-1/2}$ and
\begin{equation}
  \ket{\psi_{-n}}=\sqrt{1-\varepsilon_n^{2}}\,\ket{\tfrac12}
    +\frac{\varepsilon_n}{\sqrt2}
      \bigl(\ket{-\tfrac12}+\ket{\tfrac32}\bigr),
  \label{eq:ce-vector}
\end{equation}
and let $\varrho$ be the product state with these vectors at the sites
$-n$ and $\ket{-\tfrac12}$ at every site $x\ge0$. Then $\varrho$ lies in the
$\ell^1$ kink class \emph{including} its optional spatial first moment, since
$\bra{\psi_{-n}}S^z\ket{\psi_{-n}}=\tfrac12$ \emph{exactly} at every site: all
the deviation sums vanish identically. Yet the window charge on
$W_N=[-N,N]$ with $c_0=-1$ has, in the cyclic vector, the distribution of
$L_N=\sum_{n=1}^{N}(S^z_{-n}-\tfrac12)$, whose characteristic function is
\begin{equation}
  \varphi_N(t)=\prod_{n=1}^{N}\bigl[1-\varepsilon_n^{2}(1-\cos t)\bigr].
  \label{eq:ce-charfun}
\end{equation}
With $V_N=\sum_{n\le N}\varepsilon_n^{2}\sim\log N\to\infty$ and
$1-\cos t\ge 2t^2/\pi^2$ on $[-\pi,\pi]$, one gets
$\abs{\varphi_N(t)}\le e^{-2V_Nt^{2}/\pi^{2}}$ and hence
$\sup_k \mathbb{P}(L_N=k)\le C\,V_N^{-1/2}$: the distribution spreads without
bound. Splitting the resolvent matrix element at $\abs{k}\le K$ then gives
$\braket{\Omega_\varrho,(\Qhat_{W_N,-1}-\mathrm{i})^{-1}\Omega_\varrho}\to0$,
which is impossible for a self-adjoint limit, since every self-adjoint
$\Qhat$ has
$\operatorname{Im}\braket{\Omega,(\Qhat-\mathrm{i})^{-1}\Omega}
 =\int(\lambda^{2}+1)^{-1}\mathrm{d}\mu(\lambda)>0$. The moral is sharp: the
$\ell^1$ class controls \emph{means}, not charge fluctuations, and certainly
not implementability. A second, purely definitional point reinforces it: the
kink sector is a \emph{set} of states, whereas a GNS representation belongs to
one state, so ``the GNS representation of $\Kk_{\alpha\beta}$'' has no
canonical meaning without an additional folium choice.

\medskip\noindent\textbf{Mechanism B (on-folium; a strengthening, not a second
proof of the same statement).} This mechanism assumes strictly more --- the
smoothness hypothesis (S), the identification of the calibration parameter
with the tail density, a two-sided decorated kink reference, the corner-A
construction, and a nonscalar $V_{\theta_0}$ --- and in exchange it kills the
window-limit route \emph{on} the kink folium itself, where fluctuations are
perfectly tame. For nested windows differing by a shell $I$ of length $\ell$
deep in the $\alpha$ tail, the on-site summands commute, so
\begin{equation}
  \bigl\lVert\bigl(e^{\mathrm{i}\theta\Qhat_{W',c_0}}
    -e^{\mathrm{i}\theta\Qhat_{W,c_0}}\bigr)\Omega\bigr\rVert^{2}
  =2-2\operatorname{Re}
    \braket{\Omega,e^{-\mathrm{i}\theta Q_I^{\alpha}}\Omega},
  \label{eq:mechB}
\end{equation}
where $Q_I^{\alpha}=\sum_{x\in I}(S^z_x-\rho)$ is the tail-normal-ordered
shell charge --- the identification of the two is exactly where $s=\rho$ is
consumed. Pushing the shell deep into the tail and then letting
$\ell\to\infty$ sends the right-hand side to
$2\bigl(1-\abs{\tr(V_\theta r)}^{2}\bigr)$, and for a positive normalised
environment $\abs{\tr(Vr)}\le1$ with equality only for scalar $V$. Hence for
nonscalar $V_{\theta_0}$
\begin{equation}
  \liminf\ \bigl\lVert\bigl(e^{\mathrm{i}\theta_0\Qhat_{W'}}
    -e^{\mathrm{i}\theta_0\Qhat_{W}}\bigr)\Omega\bigr\rVert^{2}
  \ \ge\ 2\bigl(1-\abs{\tr(V_{\theta_0}r)}^{2}\bigr)>0 :
  \label{eq:mechB-bound}
\end{equation}
the window unitaries are not Cauchy, so the window charges have no
strong-resolvent limit. (At $\theta\in2\pi\Z$ there is no obstruction, exactly
as Remark \ref{ref:zero-offset} predicts.)

\medskip\noindent\emph{What survives.} Mechanism A shows the operator form
fails off every kink folium; mechanism B shows the \emph{window route} fails
on the folium as well. The surviving statements are: the finite-window
theorem (Theorem \ref{thm:fin}) and the ordered limit law (Theorem
\ref{thm:spec}), which never form such an operator; the implementer theorem
(Theorem \ref{thm:la-folium}), which reaches a self-adjoint charge on one
fixed folium by an entirely different route; and the arithmetic implication
that \emph{if} a self-adjoint $\Qhat$ is supplied and $e^{2\pi\mathrm{i}\Qhat}$
is scalar, then $\Qhat$ is pure point with spectrum in one coset of $\Z$ ---
which is a consequence of spectral calculus and proves neither premise.
\end{refutedclaim}

\provenance{M-INDEX-LA-strong}
{two mechanisms: \texttt{theory/memory-index.md} step (1.3) (the
 fluctuation counterexample, off-folium --- this alone refutes the row as
 stated and needs neither (S) nor $s=\rho$) and step (1.12) (the nonscalar
 virtual-circle obstruction, on-folium --- a \emph{strengthening}, not a
 second proof of the same statement: it additionally assumes (S), $s=\rho$,
 D1(e$'$), A2, and nonscalar $V_{\theta_0}$). The valid arithmetic
 implication is step (1.2)}
{\texttt{theory/checks/memory\_index\_check.py} IDX-C3, IDX-C8 (green exit
 $0$, \texttt{--red} exit $1$)}
{\texttt{memory-index-b-r1.md} objection 7(d) (nonexistence direction only);
 \texttt{memory-index-r2.md} objection 4; \texttt{blitz-la-folium-r1.md}
 (compatibility with Theorem \ref{thm:la-folium})}

\subsection{Where memory can live: the compact-group scope}

So far the symmetry has been one selected circle. For a general compact
on-site group one can ask which directions of the symmetry algebra are
capable of carrying a calibrated memory at all. The answer is sharp, and it is
a genuine restriction: only the central directions of the unbroken subgroup.

\begin{theorem}[Calibrated memory lives on the centre of the unbroken
subgroup]
\label{thm:scope-center}
\statusProved\
Let $G$ be a compact on-site symmetry group and $\alpha$ a vacuum of the
covariant family, with unbroken subgroup $H_\alpha$ and Lie algebra
$\mathfrak{h}_\alpha$. Then:
\begin{enumerate}
\item[(i)] the vacuum density functional
  $\xi\mapsto\omega_\alpha(q(\xi))$ is $\Ad(H_\alpha)$-invariant and
  annihilates $[\mathfrak{h}_\alpha,\mathfrak{h}_\alpha]$; with an
  $\Ad(H_\alpha)$-invariant inner product its representing vector lies in
  $\mathfrak{h}_\alpha^{H_\alpha}=\mathrm{Lie}\bigl(Z(H_\alpha)^{0}\bigr)$, and
  every nonzero tail-density jump descends through the on-site scalar kernel
  to $\mathrm{Lie}\bigl(Z_{\alpha,\mathrm{eff}}^{0}\bigr)$;
\item[(ii)] a Weyl representative $n_w$ with $\beta=n_w\cdot\alpha\neq\alpha$
  and $\Ad_{n_w^{-1}}\xi=-\xi$ supplies an opposite-density kink pair
  $(+s,-s)$ on a common unbroken circle, hence a kink sector of the
  corner-A type; and
\item[(iii)] semisimple unbroken directions and finite groups supply no
  calibrated memory: in the first case the vacuum density vanishes, so no
  positive calibration $s$ exists; in the second there is no Lie-algebra
  circle, hence no current to integrate and no real displacement calibration
  at all.
\end{enumerate}
\end{theorem}

\begin{scope}
The theorem classifies the only possible \emph{location} and \emph{lattice} of
a memory effect; it does not assert that any given compact group actually has
one. It does not claim that every $G$ has an active centre, a Weyl element
that flips a central direction, a kink, scattering channels, or relaxation. A
Weyl element supplies a calibrated pair only on a direction in its $-1$
eigenspace and only when it moves $\alpha$ to a distinct vacuum; directions
fixed by that element have equal tail densities and are spectators. The
exclusion in (iii) concerns the semisimple part of the \emph{unbroken}
algebra, not of the ambient group: an ambient semisimple group may perfectly
well exhibit memory after breaking to a subgroup with a central circle, as in
$SU(2)\to U(1)$. Finite groups may still leave exact group-valued string
endpoint operators, but that is a different observable and not an additive
memory law. Finally, this construction supplies vacuum-pair hypotheses only;
it supplies neither scattering hypotheses nor (LR).
\end{scope}

\emph{Idea of proof.} For (i), invariance of the vacuum under $H_\alpha$ gives
$d_\alpha(\Ad_h\xi)=d_\alpha(\xi)$; differentiating along $h=\exp(t\eta)$
gives $d_\alpha([\eta,\xi])=0$, so the functional kills the commutator ideal.
Compact reductivity splits
$\mathfrak{h}_\alpha=\mathfrak{z}(\mathfrak{h}_\alpha)\oplus
[\mathfrak{h}_\alpha,\mathfrak{h}_\alpha]$, and the invariance places the
representing vector in the fixed-point set, i.e. in the Lie algebra of the
connected centre. A direction acting by scalars on site has the same density
in every state, so a nonzero \emph{jump} descends to the effective centre. For
(ii), covariance gives $d_\beta(\xi)=d_\alpha(\Ad_{n_w^{-1}}\xi)=-d_\alpha(\xi)$,
which is the antisymmetric pair, and conjugation carries the same circle into
$H_\beta=n_wH_\alpha n_w^{-1}$, so the circle is common-unbroken; the
corner-A construction then supplies and superselects the kink sector. For
(iii), the wall calibration carries an explicit factor $(2s)^{-1}$ with $s>0$,
which (i) makes unavailable on semisimple directions, while a finite group has
trivial Lie algebra.

\provenance{M-SCOPE-center; D2, D9, D10, A2, M-quant-G, S-IDX-fin-G,
 S-IDX-G-label}
{\texttt{theory/compact-g-memory-scope.md} \S\S1, 3--4, i.e. steps (1.1),
 (1.3), (1.4), (1.5)}
{\texttt{theory/checks/scope\_g\_check.py} SG-C1 (finite covariance only: 24
 deterministic covariant tensors and 24 invariant density samples for a
 nonabelian $U(1)\times SU(2)$, checking zero density on all three semisimple
 generators and a permitted nonzero central density; the red mutation adds a
 highest-weight triplet component and breaks covariance)}
{\texttt{blitz-scope-g-r1.md} (0 fatal / 0 major)}

\subsection{Higher rank: the torus generalization}

When the effective centre has rank $r>1$ the escaped charge is a vector, and
the natural home for it is not $\Z^r$ but the weight lattice of the effective
central torus; $\Z^r$ appears only after a choice of primitive coordinates.

\begin{theorem}[The joint escaped charge lies in the central weight lattice]
\label{thm:torus}
\statusProved\
Let the vacuum pair share a connected torus of rank $r$ inside
$Z(H_\alpha)^0\cap Z(H_\beta)^0$; quotient its on-site scalar kernel and
choose $r$ primitive cocharacters of the resulting effective torus, so that
the free part of its character lattice is written $\Z^r$. Impose (INT)
separately on each corresponding Hermitian generator $C_j=-\mathrm{i}q(\xi_j)$,
and on every \emph{active} direction impose the antisymmetric calibration
$\omega_\alpha(C_j)=+s_j$, $\omega_\beta(C_j)=-s_j$ with $s_j>0$ (directions
with zero density jump are simply omitted from the calibrated list). Use one
fixed window $W=[a,b]$, one cut $c_0\in W$, and the same backgrounds at both
measurement times. Then:
\begin{enumerate}
\item[(i)] at finite times the joint TPM escaped-charge vector lies in the
  effective central weight lattice, hence in $\Z^r$ in the chosen primitive
  coordinates; and
\item[(ii)] under the joint-PVM strengthening of (LR) stated below, every
  ordered subsequential law $p$ on that lattice satisfies, componentwise,
  \begin{equation}
    \delta x_j=-\frac{1}{2s_j}\sum_{\nu\in\Z^r}\nu_j\,p(\nu),
    \qquad j=1,\dots,r .
    \label{eq:torus-ledger}
  \end{equation}
\end{enumerate}
The joint-PVM strengthening is: the \emph{joint} TPM laws have the (LR1)
limits on one common sequence; the \emph{joint} initial dephasing has
vanishing first-moment defect for every $C_j$, as in (LR2); and the laws are
first-moment tight with $\abs{\nu}$ replaced by $\lVert\nu\rVert_1$, as in
(LR3). The optional full-sequence weak convergence plays exactly the
convenience role it plays in Definition \ref{def:lr}.
\end{theorem}

\begin{scope}
Conclusion (i) is unconditional on the displayed finite hypotheses and needs
no relaxation assumption of any kind. Conclusion (ii) needs the joint-PVM
strengthening, and the distinction matters: applying the scalar hypothesis
(LR) separately to each circle proves the $r$ \emph{marginal} statements of
Theorem \ref{thm:spec}, but it does not by itself prove that the finer joint
dephasing defect vanishes. An individual calibrated history outcome is the
vector $\bigl(-\nu_1/(2s_1),\dots,-\nu_r/(2s_r)\bigr)$; each $\delta x_j$ is
its mean and need not lie on the lattice --- no mean quantization is claimed
(Remark \ref{ref:register-discipline}). No single scalar wall displacement is
inferred unless the event separately makes the component values compatible.
Under the smoothness hypothesis (S) at both tails, Theorem \ref{thm:density}
additionally gives $2s_j\in\Z$ for every antisymmetric active circle, but that
is a conclusion, not an added hypothesis.
\end{scope}

\emph{Idea of proof.} The generators $C_j$ lie in one abelian torus, so they
commute on site; hence all the window charges
\begin{equation}
  \Qhat^{(j)}_{W,c_0}:=2s_j\bigl(\wallX^{(j)}_W-c_0\bigr)
   =\sum_{x=a}^{b}C_{j,x}+s_j\,(a+b-1-2c_0)\,\mathbb{1}
  \label{eq:torus-charges}
\end{equation}
commute and admit a joint spectral resolution, and so do their Heisenberg
translates at each fixed time. The joint on-site spectral labels form an
affine torsor over the effective central character lattice, so every
\emph{difference} of labels evaluates on the chosen primitive cocharacters to
an element of $\Z^r$; summing over the window and adding the scalar
backgrounds changes only the common affine offset. Every nonzero atom of the
joint TPM law then has the form
\begin{equation}
  \bigl\lVert \specres_{W,t_+}(\{\lambda-\nu\})\,
              \specres_{W,t_-}(\{\lambda\})\,\Psi\bigr\rVert^{2},
  \qquad \nu\in\Z^{r},
  \label{eq:torus-atom}
\end{equation}
and the affine offset cancels because both measurements use the same window,
cut and backgrounds --- the vector version of Theorem \ref{thm:fin}(iii),
again with no cross-time commutation used. For (ii), the componentwise
analogue of the mean identity \eqref{eq:tpm-mean} plus the joint (LR2) gives
$\sum_\nu\nu_j p_W(\nu)
 =-2s_j[\omega^{+}_W(\wallX^{(j)}_W)-\omega^{-}_W(\wallX^{(j)}_W)]$, and
vector first-moment tightness --- $\lVert\nu\rVert_1$ dominates every
$\abs{\nu_j}$ --- supplies the Prokhorov and first-moment steps of Theorem
\ref{thm:spec} in each component.

\provenance{M-INDEX-torus; D13, D26, D27, M-INDEX-fin, M-INDEX-spec,
 M-IDX-density, S-IDX-G-label}
{\texttt{theory/compact-g-memory-scope.md} \S2, step (1.2), with the
 finite-window part at (1.2).(2.1)--(2.5) and the ordered-limit part at
 (1.2).(2.6)--(2.9)}
{\texttt{theory/checks/scope\_g\_check.py} SG-C2 (finite arithmetic only; no
 gate bears on the relaxation hypothesis): three commuting on-site Cartan
 charges, two random same-time joint eigenbases with noncommuting cross-time
 observables, and the full $625$-atom joint TPM law, checking normalisation
 and affine-offset cancellation into $\Z^3$; the red mutation shifts the
 final background by $\sqrt2/10$ and the lattice gate fails}
{\texttt{blitz-scope-g-r1.md} (0 fatal / 0 major; the citation for the
 on-site joint spectrum was repaired there to D26 and Theorem
 \ref{thm:fin})}

\subsection{What the section does and does not settle}

Theorem \ref{thm:fin} is the centrepiece and is unconditional: given the
integrality hypothesis on the on-site charge and a fixed window, the outcomes
of the stated two-measurement protocol are integers, for any dynamics that
acts by automorphisms. Theorem \ref{thm:spec} carries that statement to the
ordered limit at the price of one explicitly named dynamical hypothesis, (LR),
which is assumed and not derived; Theorem \ref{thm:lr1-gen} removes its first
clause from the price list, leaving the substantive content in the second and
third. Theorem \ref{thm:density} is a by-product of independent interest, and
Theorem \ref{thm:scope-center} together with Theorem \ref{thm:torus} says
where in a general compact symmetry such a memory can live and what lattice it
takes values in.

Two limits of the analysis should be kept in view. The operator form of the
finite-window statement is false (Refuted claim \ref{ref:la-strong}), so
``the total charge of the sector'' is not an available object; only the
law-level statements and the single-folium implementer (Theorem
\ref{thm:la-folium}) survive. And no model instance of (LR2)--(LR3) has been
exhibited: what is proved is a conditional whose hypothesis class is
constrained from several sides in \cref{sec:local-relaxation}. The status of
the triangle as a whole is assembled in \cref{sec:triangle-status}.
